\section{What is a Pseudorandom Number Generator?}

A pseudorandom number generator (PRNG) is an algorithm that produces a sequence of numbers that approximates the properties of truly random sequences.

Unlike true randomness, which is based on physical and inherently unpredictable processes, a PRNG is fully deterministic. This means that if the initial state (seed) is known, the entire sequence can be reproduced exactly.

\subsection{Formal Definition}

Formally, a PRNG can be described as a discrete dynamical system consisting of:

\begin{itemize}
    \item a finite state space $S$,
    \item a deterministic transition function $f : S \to S$,
    \item an output function $g : S \to \mathbb{R}$.
\end{itemize}

The generator is defined by the recurrence relation:

$$
x_{n+1} = f(x_n),
$$

with initial state (seed) $x_0 \in S$, and output sequence:

$$
u_n = g(x_n).
$$

Here, $(x_n)$ represents the internal state sequence, while $(u_n)$ represents the generated pseudorandom numbers.

\subsection{Seed (Initial Value)}

The seed $x_0$ is the initial element of the state space $S$. It uniquely determines the entire sequence generated by the PRNG.

If the same seed is used multiple times, the generator produces the same sequence, which is a key property used for reproducibility in simulations, debugging, and numerical experiments.

Different seeds lead to different trajectories in the state space, increasing variability of outputs.

\subsection{State Space and Periodicity}

Since the state space $S$ is finite and the generator is deterministic, the sequence $(x_n)$ must eventually repeat. This follows from the pigeonhole principle: after a finite number of steps, some state must reoccur, causing the sequence to enter a cycle.

Thus, every PRNG becomes eventually periodic.

The period of a PRNG is the smallest positive integer $T$ such that:

$$
x_{n+T} = x_n \quad \text{for all } n \ge n_0,
$$

for some transient length $n_0 \ge 0$.

A good PRNG is characterized by a very long period, making repetition practically unobservable in applications.

\subsection{Uniformity and Equidistribution}

Uniformity refers to how evenly the output values are distributed over a given interval.

For normalized outputs $u_n \in [0,1]$, a sequence is considered uniformly distributed if, for any subinterval $[a,b] \subset [0,1]$, we have:

$$
\lim_{N \to \infty}
\frac{1}{N} \#\{1 \le n \le N : u_n \in [a,b]\}
= b - a.
$$

This property is also known as \textit{equidistribution}.

Lack of uniformity introduces bias and can significantly distort results in simulations and statistical applications, such as Monte Carlo methods.

\subsection{Computational Randomness}

Although PRNGs are deterministic, their goal is to produce sequences that are statistically indistinguishable from truly random sequences under practical tests.

In other words, while they are not random in a physical sense, they exhibit \textit{computational randomness}, meaning that without knowledge of the seed, the sequence appears random for algorithmic and statistical purposes.

\subsection{Summary}

A pseudorandom number generator is a deterministic system defined on a finite state space that produces sequences with long periods and strong statistical properties.

A good PRNG should satisfy:
\begin{itemize}
    \item determinism (for reproducibility),
    \item long period,
    \item efficient computation,
    \item good statistical properties such as uniformity and equidistribution.
\end{itemize}

Many classical PRNGs, such as linear congruential generators, Lehmer generators, and LFSR-based systems, are specific implementations of this general framework, defined by particular choices of the transition function $f$ and output function $g$.